\section{Локальный обратимый автомат}
\label{sec:descent-ca}

Когда на каждом узле конечный алфавит, время дискретно, а обновление локально,
микроуровень~$\layerMicro$ есть \emph{клеточный автомат}:
\[
  s(x,t+1)=f\bigl(\{s(y,t)\}_{y\in N(x)}\bigr),
\]
но с физическими ограничениями из спуска: сохранение нормы, обратимость, спин~$1/2$, заряд~$\in\spaceZ$.

\begin{definition}[Закон эволюции]
\label{def:g-ca}
Пусть $\symCapitalPsi^t=\{z(x,t)\}_{x\in\carrierLattice}$ --- глобальная конфигурация (\cref{def:field-a3}).
\textbf{Закон эволюции} $g$ на~$\layerMicro$ --- локальное обратимое отображение
\[
  \symCapitalPsi^{t+1}=g(\symCapitalPsi^t),\qquad g^{-1}\circ g=\mathrm{id},
\]
при конечном регистре $Z\in\spaceZ_{\phaseRingCount}[i]$ на каждом узле.
Исторические примеры (Конвей, Rule~110) иллюстрируют класс CA, но не фиксируют~$g$;
см. приложение~\ref{app:background}.
\end{definition}

\subsection{Словарь: термины клеточного автомата и слой~\texorpdfstring{$M$}{M}}
\label{sec:ca-terminology}

В литературе по клеточным автоматам и «цифровой физике'' те же объекты часто называют иначе,
чем в QED и в гл.~\ref{ch:carrier} о двух скоростях.
Ниже --- соответствие без отождествления микро- и макроскорости.

\begin{table}[ht]
\centering
\caption{КА-терминология и наша запись на~$M$.}
\label{tab:ca-glossary}
\small
\begin{tabular}{@{}p{3.4cm}p{4.8cm}p{4.6cm}@{}}
\toprule
Термин КА & На~$M$ & Замечание \\
\midrule
ячейка, site & узел $\hv\in\carrierLattice$ & регистр $Z\in\spaceZ_{\phaseRingCount}[i]$ \\
состояние $s(x)$ & спинор $z(x)\in\spaceC^2$ & в коде --- фиксированная точка на $Z$ \\
окрестность & $N(x)$ (\cref{def:neighborhood}) & световая за один такт, не Мур \\
правило, $f$ & закон $g$ (\cref{def:g-ca}) & обратим; $g^{-1}\circ g=\mathrm{id}$ \\
поколение, generation & такт $dt=\htick$ & $t\in\spaceZ$ \\
шаг решётки & $dh=\caStepSpace$ & один hop по $\carrierLattice$ \\
\emph{скорость света} CA & $\caSignalSpeed=dh/dt=\caStepSpace/\htick$ & \textbf{тактовая} скорость сигнала \\
макроскорость & $\speedC=\kgeom \caSignalSpeed$ & мост~$\layerMacro$, CODATA (\cref{sec:two-c}) \\
световой конус & $|x-y|\le \caSignalSpeed\cdot\caStepTime\cdot dh$ & каузальность одного такта \\
алфавит & $\spaceZ_{\phaseRingCount}[i]\times Q(2^{-6})$ & конечный, не $\{0,1\}$ \\
\bottomrule
\end{tabular}
\end{table}

В учебниках по CA «скорость света'' почти всегда означает \emph{не} экспериментальную $\speedC$,
а нормировку «одна ячейка за один такт'': $\caSignalSpeed=1$ в natural units.
У нас то же: за один~$dt$ сигнал не дальше одного ребра длины~$dh$ (\cref{def:step-tick}).
Макроскопическая скорость света в вакууме~$\speedC$ появляется только после осреднения на~$\layerMacro$
и связана с~$\caSignalSpeed$ коэффициентом~$\kgeom$ (гл.~\ref{ch:carrier}, \cref{sec:two-c}).
Путать $\caSignalSpeed$ и $\speedC$ --- типичная ошибка при чтении CA-литературы рядом с физикой поля.

Закон~$g$ не подбирается под таблицу Rule~110: он \emph{замыкается} из абсолютных условий
и геометрии носителя (гл.~\ref{ch:axioms}, \ref{ch:carrier}, \ref{ch:evolution}).
Макромир~$\layerMacro$, с которого начался спуск, --- осреднение $\opBoundary$ по~$\layerMicro$;
гладкие уравнения QED и Маделунга --- его имена, не ядро автомата.

